\begin{longtable}[]{@{}
  >{\raggedright\arraybackslash}p{(\linewidth - 2\tabcolsep) * \real{0.5000}}
  >{\raggedright\arraybackslash}p{(\linewidth - 2\tabcolsep) * \real{0.5000}}@{}}
\caption{\centering طرح درس: مدل‌سازی شناختی (\lr{Cognitive Modeling)}}\tabularnewline
\toprule\noalign{}
\endfirsthead
\endhead
\bottomrule\noalign{}
\endlastfoot
\textbf{نام درس} & مدل‌سازی شناختی (درس اختیاری تخصصی) \\
\textbf{واحد} & 2 واحد: 1 واحد نظری (16 ساعت) + 1 واحد عملی (32 ساعت) =
48 ساعت \\
\textbf{شکل برگزاری} & 16 هفته × (1 ساعت درس + 2 ساعت \lr{lab} عملی) \\
\textbf{مخاطبان} & دانشجویان کارشناسی ارشد و سال آخر کارشناسی در علوم
شناختی، روان‌شناسی، علوم کامپیوتر و هوش مصنوعی \\
\textbf{پیش‌نیازها} &
روش‌های آماری در علوم شناختی؛ روش‌های آزمایشی در روان‌شناسی شناختی. فرض بر این است که دانشجو با احتمال، آمار و مقدماتی آشناست؛ آشنایی با آمار بیزی لازم نیست . \\
\textbf{مدرس} & \emph{{[}نام، اتاق، ایمیل، ساعت حضور{]}} \\
\textbf{دستیار آموزشی} & \emph{{[}نام، ایمیل، ساعت‌های \lr{lab{]}}} \\
\textbf{صفحه‌ی درس} & \emph{{[}لینک سامانه‌ی آموزشی \lr{(LMS){]}}} \\
\end{longtable}

\begin{center}\rule{0.5\linewidth}{0.5pt}\end{center}

\subsection{شرح درس}

مدل‌های شناختی توصیف‌های صریح، ریاضی یا محاسباتی، از فرایندهایی هستند که
پشت رفتار قرار دارند. در این درس یاد می‌گیریم چنین مدل‌هایی را بسازیم، روی
داده‌های آزمایشی برازش \lr{(fit)} کنیم، ارزیابی کنیم، و تفسیر کنیم که
درباره‌ی ذهن چه می‌گویند. درس را با مدل‌هایی از حوزه‌های حافظه، ادراک،
دسته‌بندی \lr{(categorization)}، تصمیم‌گیری و یادگیری، و ارتباط آن‌ها با
هوش مصنوعی تمام می‌کنیم.

\subsection{اهداف یادگیری}

\textbf{اهداف کلی.} در پایان درس، دانشجویان با موارد زیر آشنا خواهند
بود: 1. رویکردها و مدل‌های مختلف مدل‌سازی شناختی؛ 2. چگونگی طراحی یک مدل
برای حل یک مسئله‌ی شناختی؛ 3. کاربردهای مدل‌سازی در علوم شناختی و هوش
مصنوعی.

\textbf{اهداف ویژه.} دانشجویان مهارت‌های زیر را کسب می‌کنند: 1. صورت‌بندی
روش‌های مدل‌سازی روی داده‌های آزمایشی؛ 2. تفسیر نتایج پژوهشی به‌دست‌آمده با
مدل‌های شناختی.

\subsection{سرفصل‌ها}

\begin{enumerate}
\def\labelenumi{\arabic{enumi}.}
\tightlist
\item
  مبانی مدل‌سازی شناختی
\item
  روش‌های برآورد پارامتر (کمترین مربعات / میانگین مربعات خطا یا \lr{MSE}،
  بیشینه درست‌نمایی)
\item
  ارزیابی مدل‌های شناختی
\item
  مدل‌سازی بیزی
\item
  برآورد پارامتر با روش‌های داده‌محور \lr{(Monte Carlo}، نمونه‌گیری
  \lr{Metropolis\lr{--}Hastings} و \lr{Gibbs}، \lr{MCMC)}
\item
  نمونه‌هایی از مدل‌سازی در حوزه‌های مختلف علوم شناختی
\end{enumerate}

\subsection{برنامه‌ی هفتگی}

\begin{longtable}[]{@{}
  >{\raggedright\arraybackslash}p{(\linewidth - 8\tabcolsep) * \real{0.2000}}
  >{\raggedright\arraybackslash}p{(\linewidth - 8\tabcolsep) * \real{0.2000}}
  >{\raggedright\arraybackslash}p{(\linewidth - 8\tabcolsep) * \real{0.2000}}
  >{\raggedright\arraybackslash}p{(\linewidth - 8\tabcolsep) * \real{0.2000}}
  >{\raggedright\arraybackslash}p{(\linewidth - 8\tabcolsep) * \real{0.2000}}@{}}
\caption{\centering برنامه‌ی هفتگی}\tabularnewline
\toprule\noalign{}
\begin{minipage}[b]{\linewidth}\raggedright
هفته
\end{minipage} & \begin{minipage}[b]{\linewidth}\raggedright
فصل
\end{minipage} & \begin{minipage}[b]{\linewidth}\raggedright
درس نظری
\end{minipage} & \begin{minipage}[b]{\linewidth}\raggedright
\lr{lab}
\end{minipage} & \begin{minipage}[b]{\linewidth}\raggedright
موعد تحویل
\end{minipage} \\
\midrule\noalign{}
\endfirsthead
\toprule\noalign{}
\begin{minipage}[b]{\linewidth}\raggedright
هفته
\end{minipage} & \begin{minipage}[b]{\linewidth}\raggedright
فصل
\end{minipage} & \begin{minipage}[b]{\linewidth}\raggedright
درس نظری
\end{minipage} & \begin{minipage}[b]{\linewidth}\raggedright
\lr{lab}
\end{minipage} & \begin{minipage}[b]{\linewidth}\raggedright
موعد تحویل
\end{minipage} \\
\midrule\noalign{}
\endhead
\bottomrule\noalign{}
\endlastfoot
1 & 1 & مدل شناختی چیست؟ سطوح تحلیل \lr{Marr}؛ نمادگذاری &
شبیه‌سازی \lr{random walk} & \lr{---} \\
2 & 1 & کالبدشکافی یک مدل: \lr{SDT} و منحنی‌های فراموشی & \lr{SDT} و
\lr{ROC/zROC}؛ منحنی فراموشی با نویز دوجمله‌ای & \lr{A0} \\
3 & 2 & کمترین مربعات، سطح خطا و بهینه‌سازی & جستجوی شبکه‌ای \lr{(grid
search)}، \lr{Nelder\lr{--}Mead}، \lr{L-BFGS-B}؛ حالت‌های شکست & \lr{---} \\
4 & 2 & بیشینه درست‌نمایی & \lr{MLE} دوجمله‌ای؛ \(G^2\)؛ \lr{SDT} چندشرطی
& \lr{A1} \\
5 & 2 & عدم قطعیت؛ برازش برای چند شرکت‌کننده & بازه‌های اطمینان
\lr{Hessian}، \lr{profile} و \lr{bootstrap}؛ خطاهای ناشی از میانگین‌گیری
& \lr{---} \\
6 & 3 & برازش در برابر تعمیم‌پذیری؛ \lr{cross-validation} &
\lr{overfitting}؛ \lr{K-fold CV}؛ معیار تعمیم & \lr{A2} \\
7 & 3 & \lr{LRT}، \lr{AIC/BIC}؛ بازیابی پارامتر و مدل & مطالعات بازیابی؛
\lr{confusion matrix} & تشکیل گروه‌های پروژه \\
8 & 1\lr{--}3 & مرور میان‌ترم: کل مسیر مدل‌سازی & \textbf{آزمون میان‌ترم} (90
دقیقه) + مشاوره‌ی پروژه & \lr{---} \\
9 & 4 & مدل‌سازی بیزی 1: پیشین، پسین، شبکه & \lr{Beta-binomial}؛ پسین
شبکه‌ای دوبعدی برای \lr{SDT}؛ بررسی پیش‌بینی پیشین & \lr{A3}، پیشنهاد
پروژه \\
10 & 4 & مدل‌سازی بیزی 2: \lr{PPC} و \lr{Bayes factor} & \lr{PPC}؛
\lr{Bayes factor} با شبکه و \lr{Savage\lr{--}Dickey}؛ حساسیت به پیشین &
\lr{---} \\
11 & 5 & \lr{Monte Carlo} و \lr{Metropolis\lr{--}Hastings} & \lr{Metropolis}
از صفر؛ تنظیم \lr{(tuning)} & \lr{---} \\
12 & 5 & نمونه‌گیری \lr{Gibbs}؛ تشخیص همگرایی &
\lr{Gibbs }از صفر؛ \lr{R-hat }و \lr{ESS}؛ & \lr{A4} \\
13 & 5 & احتمالاتی؛ مدل‌های سلسله‌مراتبی & \lr{SDT }سلسله‌مراتبی در & گزارش
پیشرفت \\
14 & 6 & دسته‌بندی \lr{(GCM)} و تصمیم‌گیری \lr{(drift-diffusion)} &
\lr{GCM} در برابر \lr{prototype}؛ \lr{EZ-diffusion} & \lr{A5} \\
15 & 6 & یادگیری تقویتی، شناخت بیزی و هوش مصنوعی & برازش و بازیابی
\lr{Q-learning}؛ مشاهده‌گر ایده‌آل & \lr{---} \\
16 & \lr{---} & جمع‌بندی و اصول مدل‌سازی خوب & \textbf{ارائه‌ی پروژه‌ها} &
ارائه‌ها \\
دوره‌ی امتحانات & 1\lr{--}6 & \lr{---} & \textbf{آزمون پایان‌ترم} (120 دقیقه) &
گزارش + (پیش از آزمون \\
\end{longtable}

\subsection{ارزشیابی (از 20)}

\begin{longtable}[]{@{}ll@{}}
\caption{\centering ارزشیابی (از 20)}\tabularnewline
\toprule\noalign{}
بخش & نمره \\
\midrule\noalign{}
\endfirsthead
\toprule\noalign{}
بخش & نمره \\
\midrule\noalign{}
\endhead
\bottomrule\noalign{}
\endlastfoot
تأیید تمرین‌های \lr{lab} (تمرین‌های هفتگی سر کلاس) & 2.0 \\
تکالیف \lr{A1 }تا \lr{A5 (}هرام 1.0؛ \lr{A0 }الزامی ولی بدون نمره &
5.0 \\
کوییزهای سر کلاس (حدود 10 کوییز؛ 2 نمره‌ی پایین حذف می‌شود) & 1.0 \\
آزمون میان‌ترم (هفته‌ی 8) & 3.0 \\
پروژه‌ی پایانی: پیشنهاد 1.0 · ارائه 2.0 · گزارش + 3.0 & 6.0 \\
آزمون پایان‌ترم (جامع، با تأکید بر فصل‌های 4 تا 6) & 3.0 \\
\textbf{جمع} & \textbf{20.0} \\
\end{longtable}

\subsection{قوانین (خلاصه)}

\begin{itemize}
\tightlist
\item
  \textbf{تأخیر در تحویل:} 10\% کسر نمره برای هر روز، حداکثر 3 روز.
\item
  \textbf{ابزارهای هوش مصنوعی:} برای رفع اشکال \lr{(debugging)}، توضیح
  مفاهیم و ویرایش زبانی مجاز است. نحوه‌ی استفاده را بنویسید؛ باید بتوانید
  همه‌ی آن‌چه تحویل می‌دهید را توضیح دهید.
\item
  \textbf{آزمون‌ها:} کتاب‌بسته، برگه‌ی فرمول داده می‌شود، ماشین‌حساب
  غیربرنامه‌پذیر مجاز است.
\item
  \textbf{حضور:} بیشترِ یادگیری در \lr{lab} اتفاق می‌افتد؛ تأیید تمرین‌های
  \lr{lab} به حضور (یا یک جایگزین توافق‌شده) نیاز دارد.
\item
  \textbf{دسترس‌پذیری:} دانشجویانی که به شرایط ویژه نیاز دارند در دو
  هفته‌ی اول با مدرس تماس بگیرند.
\end{itemize}

\subsection{منابع}

(2018). \emph{\lr{Computational modeling of cognition and behavior.}}
\lr{Cambridge University Press.} \textbf{(کتاب اصلی، هفته‌های 1 تا 8)} -
\lr{Lee, M. D.,} \lr{\&} \lr{Wagenmakers, E.-J.} (2014). \emph{\lr{Bayesian
cognitive modeling: A practical course.}} \lr{Cambridge University
Press.} \textbf{(کتاب اصلی، هفته‌های 9 تا 14)} - \lr{Kruschke, J. K.}
(2014). \emph{\lr{Doing Bayesian data analysis: A tutorial with R, JAGS,
and Stan}} \lr{(2nd ed.). Academic Press.} - \lr{Anderson, B.} (2014).
\emph{\lr{Computational neuroscience and cognitive modelling: A
student's introduction to methods and procedures.}} \lr{SAGE.} -
\lr{Verguts, T.} (2022). \emph{\lr{Introduction to modeling cognitive
processes.}} \lr{MIT Press.} - قدوسیان، ع. (1401). \emph{مبانی روش‌های
مستقیم در بهینه‌سازی غیرخطی.} - عظیمی‌فر، ز. و احمدی، ا. (1396).
\emph{شناسایی آماری الگو: یادگیری نظری مقدماتی}، جلد 1.
